\section{Reproducibility and epistemic status}
\label{app:reproducibility}

The source archive contains a standard-library Python supplement.  From the
manuscript directory, run
\begin{verbatim}
python3 supplement/representation_certificate.py \
  --output supplement/representation_certificate.json
python3 -m unittest -v supplement/test_representation_certificate.py
\end{verbatim}
Thirteen tests should pass, and the JSON field \verb|all_checks_pass| should
be true.  The certificate uses integers and \verb|fractions.Fraction| only.
It records exact profile composition, uniform and collapsing finite gaps, a
hidden sector, and the noncommuting defect array.

The trusted mathematical base of Theorem~\ref{thm:recovery-spectrum} is the
positive spectral theorem and the analytic proof in Appendix
\ref{app:spectral}.  The code is not a proof of that theorem.  The OS
corollary is conditional on the stated reconstruction and density premises.
All implications asserting the required four-dimensional Yang--Mills
continuum construction or uniform rigidity estimate are marked open.

The bibliography was checked against publisher, DOI, arXiv, and official
Clay records on 21 July 2026.  The search ledger, closest-result matrix,
claim--source matrix, proof graph, countermodel ledger, and scope audit are
included with the source.  A subscription-database priority search and author
metadata remain external pre-submission tasks.
